% =====================================================================
%  Chapter 21 — The Proof in Natural Language, 5
% =====================================================================

\chapter{The Lift Chain and the Multiplicity}
\label{ch:chain2}

We now reach the device that makes the method possible: the iterated theta
lift.

\section{The statement of the lift}Fix a CM field $E$, a positive definite Hermitian form $A$ of rank $n$ over
$E$, and a cuspidal representation $\pi$ of the unitary group
$\mathrm{U}(E,n)$. The theta series $\theta_A$ is a modular form of weight $n$
on the CM lattice.

\begin{quote}
\textbf{Proposition L (theta lift).} Let $\mu$ be a Hecke character of
$\mathrm{U}(E,n)$ and let $t$ be a parameter determining $A$. If the weight of
$\pi$ matches the weight of $\theta_A$, there is at most one representation
$\pi'$ of $\mathrm{U}(E,n')$ with $n' > n$ whose theta lift is $\pi \boxplus
\mu[t]$.
\end{quote}

The weight-matching condition is an inequality: for the lift from rank $n$ to
rank $d$ to exist, one needs
\begin{equation}
d + m_W(w) \;\le\; 2\,\lambda(w)
\label{eq:weightineq2}
\end{equation}
at every archimedean place $w$, where $\lambda$ is the weight of $\pi$ and
$m_W$ the weight attached to the form. This is the condition stated in
\code{UpperTheory.TP.TLiftK}, and it is the reason the lift exists only in
certain directions.

The formal counterpart is \code{UpperTheory.TP.tLift\_of\_tools\_k}, assembled
in \code{.../ThetaKernels/TPGHolds.lean} from roughly twenty separately proved
theorems.

\section{Iterating}One lift is not enough. Apply it again to the output, with a new form and a
new character:

\begin{quote}
\textbf{Proposition M (the chain).} There is a chain
\[
U(X_0) \longrightarrow U(X_1) \longrightarrow \cdots \longrightarrow U(X_m),
\]
at each step a theta lift $\pi \mapsto \pi \boxplus \mu_j[t]$, such that each
step preserves cuspidality, the type conditions and the non-vanishing of the
period, and such that the composition of the chain is a single lift with
multiplied kernel and multiplied weight.
\end{quote}

The last clause is the load-bearing one, and it is the formal statement of
the remark's claim that the chain ``wins multiplicity $m$''.

\begin{seamnote}{The composition law}
{\footnotesize
\begin{verbatim}
def KerComp : Prop :=
forall (R : Type) [CommSemiring R] (i k m : Type) [Fintype
    i] [Fintype k] [Fintype m]
(w1 : i -> R) (w2 : k -> R) (K1 : i -> k -> R) (K2 : k -> m
      -> R) (f : i -> R),
kerOp w2 K2 (kerOp w1 K1 f) = kerOp w1 (kerMul w2 K1 K2) f
\end{verbatim}
}
\end{seamnote}

Read it: the composition of two lifts equals a single lift whose kernel is
the product of the two kernels. The chain is therefore not a new object. It is
one lift, with a multiplied kernel, and the kernel's size is what has grown.

\section{Why that is the substitute for multiplicity}Return to the function-field method of Chapter~\ref{ch:step1}. There, the
multiplicity came from the independence of a geometric direction and an
arithmetic direction, and the dimension count multiplied the two. Here there
is no second direction, and the theta lift takes its place: the kernel of the
composed lift is the product of the kernels, so the space has grown by the
product of the growth factors.

\begin{quote}
\textbf{Proposition N.} If a single lift multiplies the relevant space by a
factor $\lambda$, then a chain of length $m$ multiplies it by $\lambda^m$.
\end{quote}

This is the statement that makes the method work, and it is the statement that
must be paid for. The remark's own phrasing of the length is the one-line
summary of the whole argument:

\begin{quote}
\small
``for the von Mangoldt sum of length $X$, Claude's suggested $\sqrt{\ }$
length of chain of theta lifts, in view of each mimicking multiplicity, was
$X$. (And yet, what else could it be?)''
\end{quote}

\section{The length $\sqrt{X}$, and why it is not obvious}A chain of length $m$ multiplies the space by $\lambda^m$, and to dominate a
sum of length $X$ one needs $\lambda^m \ge X$. For a fixed $\lambda$, that
gives $m \ge \log X / \log \lambda$ --- a logarithm, not a square root.

A square root arises if $\lambda$ itself is allowed to grow with $X$. If the
per-lift growth is $X^{\alpha}$ with $\alpha$ fixed, then a chain of length
$m$ gives $X^{\alpha m}$, and $m \ge 1/\alpha$ suffices --- a constant. If
the per-lift growth is a fixed multiple of $X$, then a chain of length $m$
gives $X^m$, and one lift would do.

So the square root is not what the multiplicativity alone predicts. It is what
one gets in an intermediate regime, and the author's parenthesis --- ``(And
yet, what else could it be?)'' --- records that he cannot account for it.

\begin{quote}
\textbf{This is the honest reading.} Either the per-lift growth factor is
itself a growing quantity, in which case the chain length is not determined by
the multiplicativity alone; or the square root enters through some other
mechanism that the remark does not describe. The development contains the
machinery that would settle it --- the budget arithmetic of
Chapter~\ref{ch:budget} --- but the specific claim does not appear in the
remark, and a reader has to find it in the code.
\end{quote}

\section{What each step must preserve}For the chain to be usable, three things must hold at each stage.

\begin{quote}
\textbf{Proposition O.} For each $j < m$:
\begin{enumerate}
\item the type condition is preserved --- the lifted object remains of the
required type;
\item the lifted representation is irreducible, hence cuspidal;
\item the Bessel period at that stage does not vanish.
\end{enumerate}
\end{quote}

These are the descent-tower conditions, and the theorem that establishes them
is \code{Conclusion1} in
\code{FunctionModelRallisFromInputs.lean}: for every admissible parameter
choice satisfying the global hypotheses \code{Hyp3} and the character
conditions \code{IsMu}, the descent-tower conditions hold.

Item 2 deserves a word. A theta lift of a cusp form is irreducible; a theta
lift of an Eisenstein series is not, and its period is a product of
$L$-values rather than a new $L$-function. The whole argument is about the
$L$-function of the zeta function and the Weil expression built from its
zeros, so it needs the cuspidal case throughout.

\section{The price}Each lift carries a denominator. Formally:

\begin{quote}
\textbf{Proposition P.} There exist an integral model $S$ and non-negative
integers $D_j$ with the property that the output of the chain of length $m$
lies in $S$ up to the factor
\[
\prod_{j < m} D_j.
\]
\end{quote}

This is \code{ChainInt}, and it is the formal statement of ``the cost of each
lift''. The product is what the argument has to pay, and the gain is the
product of the growth factors.

\begin{keytakeaways}
\item A theta lift is well defined under a weight inequality, and a chain of
them composes into a single lift with multiplied kernel --- the formal
statement of the claim that the chain wins multiplicity.
\item Each step must preserve cuspidality, the type conditions and the
non-vanishing of the period; these are the descent-tower conditions.
\item Each step also carries a denominator, and the chain's total denominator
is the product. Gain and price are both multiplicative.
\item The proposed chain length $\sqrt{X}$ does not follow from the
multiplicativity alone; the remark does not account for it, and that is
flagged here rather than smoothed over.
\end{keytakeaways}
